\section{Môi trường thực nghiệm}

\subsection{Môi trường huấn luyện}

Thực nghiệm mô hình được thực hiện trên hai môi trường. Các thực nghiệm sơ bộ
(pilot) để kiểm chứng đúng đắn của từng script --- thưa hóa, cắt tỉa, phục hồi,
chưng cất --- chạy trên máy trạm Windows chỉ dùng CPU với tập COCO128 (128 ảnh,
80 lớp; dùng chung cho train và validation). Pipeline đầy đủ trên tập
Construction-PPE chạy trên GPU của Google Colab qua notebook
\texttt{yolo26\_ppe\_train\_pipeline.ipynb}, vì QAT với \texttt{pytorch-quantization}
bắt buộc có CUDA.

\begin{table}[H]
\centering
\caption{Môi trường thực nghiệm}
\label{tab:experiment-devices}
\begin{tabularx}{\textwidth}{L{3.4cm}Y}
\toprule
\textbf{Môi trường} & \textbf{Cấu hình}\\
\midrule
Pilot (CPU) & Windows 11, Python 3.12, PyTorch CPU, Ultralytics (bản mở rộng của đồ án); COCO128; ảnh 320 hoặc 640, batch 8, \texttt{workers=0}\\
Huấn luyện đầy đủ & Google Colab GPU, đầu ra lưu trên Google Drive; Construction-PPE; ảnh 640, batch 16\\
Jetson Nano & Jetson Nano 4GB, JetPack 4.6, power mode MAXN (10 W), quạt tản nhiệt chủ động\\
Jetson Orin Nano & Jetson Orin Nano 8GB, JetPack 6, power mode 15 W, quạt tản nhiệt chủ động\\
Cloud & AWS EC2 Linux, Docker Compose (Mục 5.1)\\
\bottomrule
\end{tabularx}
\end{table}

\subsection{Quy trình đo trên thiết bị}

Mỗi engine được đo với ảnh $640\times640$, batch 1, sau 50 lần suy luận khởi động
(warm-up), trên 500 lần suy luận liên tiếp. Latency báo cáo gồm trung vị và p95;
FPS tính trên cùng khoảng thời gian. Mức sử dụng GPU, bộ nhớ, công suất và nhiệt
độ được thu bằng \texttt{tegrastats} trong bài chạy 30 phút để quan sát trạng
thái ổn định, không chỉ vài giây đầu. Power mode và xung nhịp được cố định bằng
\texttt{nvpmodel} và \texttt{jetson\_clocks} trước mỗi bài đo.
